\documentclass[10pt, xcolor={dvipsnames}]{beamer}

%===========================================================
% 1. PHẦN KHAI BÁO HỆ THỐNG
%===========================================================
\usepackage[utf8]{inputenc}
\usepackage[T5]{fontenc}
\usepackage[vietnamese]{babel}
\usepackage{listings}
\usepackage[most]{tcolorbox}

\definecolor{mainblue}{RGB}{0, 102, 104}
\definecolor{lightorange}{RGB}{255, 140, 0}
\usetheme{Madrid}
\setbeamercolor{palette primary}{bg=mainblue, fg=white}

\lstset{
    language=C++,
    basicstyle=\ttfamily\tiny,
    keywordstyle=\color{blue}\bfseries,
    commentstyle=\color{gray}\itshape,
    backgroundcolor=\color{gray!5},
    frame=single,
    breaklines=true
}

\begin{document}

%--- BÀI TOÁN 1: DFS ---
\begin{frame}{Bài 1: Chiến dịch giải cứu Voucher Buffet}
    \begin{tcolorbox}[colback=mainblue!5, colframe=mainblue, title=\textbf{Cốt truyện phiêu lưu}]
        \scriptsize 
        Nhóm \textbf{Tung Tung Tung Sahur} cần kiểm tra xem có đường đi từ vị trí $S$ đến nơi cất giấu Voucher $T$ trên \textbf{Đảo Tmath} hay không để được \textbf{Thầy Toàn} khao ăn Buffet.
    \end{tcolorbox}

    \begin{columns}[T]
        \begin{column}{0.55\textwidth}
            \begin{block}{\scriptsize Giải thích Input (Nhập theo thứ tự)}
                \tiny 
                \textbf{(1)} Dòng 1: \texttt{N M S T} (Số điểm, số đường, điểm đầu, điểm đích)\\
                \textbf{(2)} M dòng sau: \texttt{u v} (Có đường nối u và v)
            \end{block}
        \end{column}

        \begin{column}{0.4\textwidth}
            \begin{tcolorbox}[colback=white, colframe=mainblue, title=\tiny Sample Test, arc=0mm]
                \tiny
                \textbf{Input:}\\
                \texttt{4 3 1 4}\\
                \texttt{1 2}\\
                \texttt{2 3}\\
                \texttt{3 4}\\
                \vspace{0.1cm}
                \textbf{Output:}\\
                \texttt{YES}
            \end{tcolorbox}
        \end{column}
    \end{columns}
\end{frame}

%--- CODE BÀI 1 ---
\begin{frame}[fragile]{Giải thuật Bài 1: Thuật toán loang DFS}
\begin{lstlisting}
#include <bits/stdc++.h>
using namespace std;
vector<int> adj[1001]; 
bool vis[1001];

void dfs(int u) {
    vis[u] = true;
    for (int v : adj[u]) if (!vis[v]) dfs(v);
}

int main() {
    int n, m, s, t; cin >> n >> m >> s >> t;
    while (m--) {
        int u, v; cin >> u >> v;
        adj[u].push_back(v); adj[v].push_back(u);
    }
    dfs(s); 
    if (vis[t]) cout << "YES"; else cout << "NO";
    return 0;
}
\end{lstlisting}
\end{frame}

%--- BÀI TOÁN 2: DIJKSTRA CƠ BẢN ---
\begin{frame}{Bài 2: Cuộc đua đến nhà hàng Buffet}
    \begin{tcolorbox}[colback=lightorange!5, colframe=lightorange, title=\textbf{Cốt truyện}]
        \scriptsize 
        Sau khi có Voucher, nhóm cần tìm đường \textbf{ngắn nhất} từ $S$ đến nhà hàng $T$ để kịp giờ ăn. Mỗi con đường sẽ tốn một khoảng thời gian $w$.
    \end{tcolorbox}

    \begin{columns}[T]
        \begin{column}{0.55\textwidth}
            \begin{block}{\scriptsize Giải thích Input (Nhập theo thứ tự)}
                \tiny 
                \textbf{(1)} Dòng 1: \texttt{N M S T}\\
                \textbf{(2)} M dòng sau: \texttt{u v w} ($w$ là thời gian đi từ u đến v)
            \end{block}
            \begin{alertblock}{\scriptsize Gợi ý cách giải}
                \tiny Dùng mảng lưu khoảng cách. Mỗi bước chọn địa điểm chưa thăm có khoảng cách nhỏ nhất để tối ưu các điểm lân cận.
            \end{alertblock}
        \end{column}

        \begin{column}{0.4\textwidth}
            \begin{tcolorbox}[colback=white, colframe=lightorange, title=\tiny Sample Test, arc=0mm]
                \tiny
                \textbf{Input:}\\
                \texttt{3 2 1 3}\\
                \texttt{1 2 5}\\
                \texttt{2 3 10}\\
                \vspace{0.1cm}
                \textbf{Output:}\\
                \texttt{15}
            \end{tcolorbox}
        \end{column}
    \end{columns}
\end{frame}

%--- CODE BÀI 2 ---
\begin{frame}[fragile]{Giải thuật Bài 2: Dijkstra cơ bản (Dễ hiểu)}
\begin{lstlisting}
#include <bits/stdc++.h>
using namespace std;
int n, m, s, t;
int d[1001], vis[1001], a[1001][1001];

int main() {
    cin >> n >> m >> s >> t;
    // Khoi tao ma tran ke va mang khoang cach
    for(int i=1; i<=n; i++) {
        d[i] = 1e9; // So rat lon
        for(int j=1; j<=n; j++) a[i][j] = 1e9;
    }
    while(m--) {
        int u, v, w; cin >> u >> v >> w;
        a[u][v] = a[v][u] = w;
    }
    d[s] = 0;
    for(int i=1; i<=n; i++) {
        int u = 0;
        // Tim dinh u chua tham co d[u] nho nhat
        for(int j=1; j<=n; j++)
            if(!vis[j] && (u==0 || d[j] < d[u])) u = j;
        
        if(d[u] == 1e9) break;
        vis[u] = 1;
        // Toi uu khoang cach den cac dinh kề v
        for(int v=1; v<=n; v++)
            if(d[u] + a[u][v] < d[v]) d[v] = d[u] + a[u][v];
    }
    if(d[t] == 1e9) cout << -1; else cout << d[t];
    return 0;
}
\end{lstlisting}
\end{frame}

\end{document}